\documentclass[12pt,a4paper]{article}

% ---------------------------------------------------------------------------
% Relatorio de Entrega Final de Projeto Scrum
% Compilacao: pdflatex no Overleaf
% Edite os campos entre colchetes conforme os dados reais do projeto.
% ---------------------------------------------------------------------------

\usepackage[utf8]{inputenc}
\usepackage[T1]{fontenc}
\usepackage[brazil]{babel}
\usepackage{lmodern}
\usepackage{geometry}
\usepackage{graphicx}
\usepackage{xcolor}
\usepackage{array}
\usepackage{booktabs}
\usepackage{longtable}
\usepackage{tabularx}
\usepackage{multirow}
\usepackage{hyperref}
\usepackage{titlesec}
\usepackage{fancyhdr}
\usepackage[most]{tcolorbox}
\usepackage{tikz}
\usepackage{pgfplots}
\usepackage{enumitem}
\usepackage{float}
\usepackage{setspace}

\usetikzlibrary{arrows.meta, positioning, shapes.geometric}
\pgfplotsset{compat=1.18}

% ---------------------------------------------------------------------------
% Configuracoes visuais
% ---------------------------------------------------------------------------
\geometry{
  a4paper,
  left=2.5cm,
  right=2.5cm,
  top=2.5cm,
  bottom=2.5cm
}

\definecolor{PrimaryBlue}{HTML}{1F4E79}
\definecolor{DarkBlue}{HTML}{17365D}
\definecolor{SoftBlue}{HTML}{EAF2F8}
\definecolor{SoftGray}{HTML}{F4F6F7}
\definecolor{MediumGray}{HTML}{5D6D7E}
\definecolor{SuccessGreen}{HTML}{2E7D32}
\definecolor{SoftGreen}{HTML}{E8F5E9}
\definecolor{AlertGold}{HTML}{F9E79F}

\hypersetup{
  colorlinks=true,
  linkcolor=PrimaryBlue,
  urlcolor=PrimaryBlue,
  citecolor=PrimaryBlue,
  pdftitle={Relatorio de Entrega Final de Projeto Scrum},
  pdfauthor={[Nome/Papel]}
}

\setstretch{1.18}
\setlength{\parskip}{0.45em}
\setlength{\parindent}{0pt}

\pagestyle{fancy}
\fancyhf{}
\lhead{\textcolor{MediumGray}{Relatorio Final Scrum}}
\rhead{\textcolor{MediumGray}{[Nome do Projeto]}}
\cfoot{\textcolor{MediumGray}{\thepage}}
\renewcommand{\headrulewidth}{0.4pt}
\renewcommand{\headrule}{\hbox to\headwidth{\color{PrimaryBlue}\leaders\hrule height \headrulewidth\hfill}}

\titleformat{\section}
  {\Large\bfseries\color{DarkBlue}}
  {\thesection.}
  {0.6em}
  {}

\titleformat{\subsection}
  {\large\bfseries\color{PrimaryBlue}}
  {\thesubsection}
  {0.6em}
  {}

\renewcommand{\arraystretch}{1.35}

\newcolumntype{Y}{>{\raggedright\arraybackslash}X}
\newcolumntype{L}[1]{>{\raggedright\arraybackslash}p{#1}}
\newcolumntype{C}[1]{>{\centering\arraybackslash}p{#1}}

\tcbset{
  enhanced,
  boxrule=0.6pt,
  arc=2mm,
  left=3mm,
  right=3mm,
  top=2mm,
  bottom=2mm
}

\newtcolorbox{infobox}[1]{
  colback=SoftBlue,
  colframe=PrimaryBlue,
  title=\textbf{#1},
  coltitle=white,
  colbacktitle=PrimaryBlue,
  fonttitle=\bfseries
}

\newtcolorbox{successbox}[1]{
  colback=SoftGreen,
  colframe=SuccessGreen,
  title=\textbf{#1},
  coltitle=white,
  colbacktitle=SuccessGreen,
  fonttitle=\bfseries
}

\newtcolorbox{notebox}[1]{
  colback=SoftGray,
  colframe=MediumGray,
  title=\textbf{#1},
  coltitle=white,
  colbacktitle=MediumGray,
  fonttitle=\bfseries
}

\newcommand{\statusdone}{\textcolor{SuccessGreen}{\textbf{Concluido}}}
\newcommand{\statusreplaced}{\textcolor{PrimaryBlue}{\textbf{Substituido}}}
\newcommand{\placeholder}[1]{\textcolor{MediumGray}{[#1]}}

% ---------------------------------------------------------------------------
% Documento
% ---------------------------------------------------------------------------
\begin{document}

% ---------------------------------------------------------------------------
% Capa
% ---------------------------------------------------------------------------
\begin{titlepage}
  \thispagestyle{empty}
  \begin{tikzpicture}[remember picture, overlay]
    \fill[DarkBlue] (current page.north west) rectangle ([yshift=-5.5cm]current page.north east);
    \fill[PrimaryBlue] ([yshift=-5.5cm]current page.north west) rectangle ([yshift=-5.75cm]current page.north east);
  \end{tikzpicture}

  \vspace*{1.2cm}
  {\color{white}\Huge\bfseries Relatorio de Entrega Final de Projeto}\\[0.35cm]
  {\color{white}\LARGE\bfseries Scrum}\\[1.6cm]

  \begin{tcolorbox}[
    colback=white,
    colframe=PrimaryBlue,
    width=\textwidth,
    boxrule=0.8pt,
    arc=3mm
  ]
    \Large\textbf{Dados do Documento}\\[0.4cm]
    \begin{tabularx}{\textwidth}{@{}L{5cm}Y@{}}
      \textbf{Nome do Projeto:} & \placeholder{Nome do Projeto} \\
      \textbf{Periodo de Execucao:} & \placeholder{Mes/Ano} a \placeholder{Mes/Ano} \\
      \textbf{Duracao:} & 6 meses \\
      \textbf{Data de Emissao:} & \placeholder{Data} \\
      \textbf{Responsavel:} & \placeholder{Nome/Papel, ex: Scrum Master / Product Owner} \\
    \end{tabularx}
  \end{tcolorbox}

  \vfill
  \begin{flushright}
    \textcolor{MediumGray}{Documento academico/corporativo para consolidacao das entregas, metricas, retrospectiva e encerramento do ciclo principal do projeto.}
  \end{flushright}
\end{titlepage}

\tableofcontents
\newpage

% ---------------------------------------------------------------------------
\section{Executivo e Sumario do Projeto}

Este relatorio consolida as entregas realizadas ao longo de seis meses de execucao do projeto utilizando o framework Scrum. O documento apresenta a evolucao do escopo, as principais entregas, os indicadores de performance do time, os marcos do cronograma, as licoes aprendidas e os proximos passos recomendados para sustentacao e evolucao continua.

\begin{infobox}{Resumo Executivo do Projeto}
  \begin{tabularx}{\textwidth}{@{}L{5cm}Y@{}}
    \textbf{Nome do Projeto:} & \placeholder{Nome do Projeto} \\
    \textbf{Visao do Produto:} & \placeholder{Ex: Desenvolver uma plataforma web para automatizar a gestao de agentes de IA para empresas parceiras} \\
    \textbf{Data de Inicio:} & \placeholder{DD/MM/AAAA} \\
    \textbf{Data de Termino:} & \placeholder{DD/MM/AAAA} \\
    \textbf{Total de Sprints Realizadas:} & \placeholder{Ex: 12 sprints de 2 semanas} \\
    \textbf{Responsavel:} & \placeholder{Nome/Papel} \\
  \end{tabularx}
\end{infobox}

O proposito do projeto foi entregar uma solucao capaz de gerar valor direto ao cliente final por meio da digitalizacao de processos, melhoria da experiencia de uso e aumento da eficiencia operacional. Ao final do ciclo principal, o produto apresenta um conjunto de funcionalidades validado em incrementos sucessivos, com adaptacoes realizadas conforme feedbacks dos stakeholders e prioridades do negocio.

% ---------------------------------------------------------------------------
\section{Escopo Entregue vs. Objetivos do Produto}

A tabela a seguir apresenta a relacao entre os objetivos planejados, o status final e o valor gerado para o produto. Os itens podem ser ajustados conforme o backlog real do projeto.

\begin{table}[H]
  \centering
  \caption{Escopo entregue e valor gerado}
  \begin{tabularx}{\textwidth}{@{}Y C{3cm} Y@{}}
    \toprule
    \textbf{Objetivo do Produto / Epico / Feature} & \textbf{Status} & \textbf{Impacto / Valor Gerado} \\
    \midrule
    Integracao de APIs de IA & \statusdone & Permite automacao em tempo real para o cliente final. \\
    Dashboard de Metricas & \statusdone & Oferece visualizacao clara dos dados para apoio a tomada de decisao. \\
    Modulo de Relatorios Avancados & \statusreplaced & Alterado no backlog por prioridade de seguranca durante a Sprint 8. \\
    Autenticacao e Seguranca & \statusdone & Aumenta a protecao de acesso e reduz riscos operacionais. \\
    Melhorias de UI/UX & \statusdone & Melhora a experiencia dos usuarios finais. \\
    \bottomrule
  \end{tabularx}
\end{table}

\begin{notebox}{Nota do Product Owner}
  \placeholder{Inserir breves consideracoes sobre as mudancas de escopo ao longo dos seis meses, destacando a adaptabilidade do framework Scrum para responder as necessidades do mercado, feedbacks dos usuarios e prioridades do negocio.}
\end{notebox}

% ---------------------------------------------------------------------------
\section{Metricas de Performance do Time -- Agile KPIs}

\subsection{Entrega e Produtividade}

\begin{infobox}{Indicadores de Entrega}
  \begin{tabularx}{\textwidth}{@{}L{6cm}Y@{}}
    \textbf{Total de Story Points Planejados:} & \placeholder{X pontos} \\
    \textbf{Total de Story Points Entregues -- Done:} & \placeholder{Y pontos} \\
    \textbf{Taxa de Entrega -- Say-Do Ratio:} & \placeholder{X\%} \\
    \textbf{Velocidade Media do Time -- Velocity:} & \placeholder{X pontos por Sprint} \\
  \end{tabularx}
\end{infobox}

\begin{table}[H]
  \centering
  \caption{Metricas de produtividade}
  \begin{tabularx}{\textwidth}{@{}L{6.2cm}C{3.2cm}Y@{}}
    \toprule
    \textbf{Metrica} & \textbf{Valor} & \textbf{Interpretacao} \\
    \midrule
    Story Points Planejados & \placeholder{X} & Volume total estimado para o ciclo de seis meses. \\
    Story Points Entregues -- Done & \placeholder{Y} & Trabalho efetivamente aceito conforme Definition of Done. \\
    Say-Do Ratio & \placeholder{X\%} & Relação entre compromisso assumido e entrega concluida. \\
    Velocity Media & \placeholder{X pontos/Sprint} & Media de entrega por sprint, usada para previsibilidade futura. \\
    \bottomrule
  \end{tabularx}
\end{table}

\begin{figure}[H]
  \centering
  \begin{tikzpicture}
    \begin{axis}[
      ybar,
      width=\textwidth,
      height=8cm,
      bar width=12pt,
      ymin=0,
      ymax=45,
      ylabel={Story Points Entregues},
      xlabel={Sprint},
      title={Velocidade do Time por Sprint},
      symbolic x coords={S1,S2,S3,S4,S5,S6,S7,S8,S9,S10,S11,S12},
      xtick=data,
      x tick label style={rotate=45, anchor=east},
      nodes near coords,
      nodes near coords align={vertical},
      grid=major,
      major grid style={draw=SoftGray},
      axis line style={draw=MediumGray},
      tick style={draw=MediumGray},
      every axis plot/.append style={fill=PrimaryBlue, draw=DarkBlue}
    ]
      \addplot coordinates {
        (S1,20) (S2,24) (S3,28) (S4,25) (S5,30) (S6,32)
        (S7,31) (S8,35) (S9,34) (S10,36) (S11,38) (S12,40)
      };
    \end{axis}
  \end{tikzpicture}
  \caption{Valores ficticios editaveis para demonstracao da evolucao da velocidade.}
\end{figure}

\subsection{Qualidade e Bugs}

\begin{table}[H]
  \centering
  \caption{Indicadores de qualidade}
  \begin{tabularx}{\textwidth}{@{}L{6.5cm}Y@{}}
    \toprule
    \textbf{Indicador} & \textbf{Resultado} \\
    \midrule
    Debito Tecnico Mapeado & \placeholder{X\% do backlog total ou X cards} \\
    Bugs reportados pos-release em producao & \placeholder{X incidentes} \\
    SLA de correcao & \placeholder{Cumprido / Parcial / Nao cumprido} \\
    Incidentes criticos & \placeholder{X} \\
    Testes executados & \placeholder{X casos de teste} \\
    Taxa de aprovacao dos testes & \placeholder{X\%} \\
    \bottomrule
  \end{tabularx}
\end{table}

Os indicadores de qualidade demonstram a estabilidade do incremento entregue e ajudam a orientar a etapa de sustentacao do produto. A leitura conjunta de bugs, incidentes, testes e debito tecnico permite priorizar melhorias sem comprometer o valor ja disponibilizado ao usuario final.

% ---------------------------------------------------------------------------
\section{Cronograma de Entregas Importantes -- Milestones}

\begin{figure}[H]
  \centering
  \begin{tikzpicture}[
    timeline/.style={draw=PrimaryBlue, line width=1.2pt},
    milestone/.style={circle, draw=PrimaryBlue, fill=SoftBlue, minimum size=8mm, font=\bfseries\small},
    labelbox/.style={rectangle, rounded corners=2mm, draw=MediumGray, fill=SoftGray, text width=4.2cm, align=left, font=\small}
  ]
    \draw[timeline] (0,0) -- (14,0);

    \foreach \x/\num/\mes/\texto in {
      0/1/Mes 1/{Sprints 1--2\\Ambiente, arquitetura base e MVP da interface},
      2.8/2/Mes 2/{Sprints 3--4\\Primeiro modulo core em homologacao},
      5.6/3/Mes 3/{Sprints 5--6\\Primeira release em producao e beta interna},
      8.4/4/Mes 4/{Sprints 7--8\\Feedbacks e performance},
      11.2/5/Mes 5/{Sprints 9--10\\Seguranca e relatorios},
      14/6/Mes 6/{Sprints 11--12\\Carga, UI/UX e rollout final}
    }{
      \node[milestone] at (\x,0) {\num};
      \node[font=\bfseries\small, text=DarkBlue] at (\x,0.75) {\mes};
      \node[labelbox, anchor=north] at (\x,-0.75) {\texto};
    }
  \end{tikzpicture}
  \caption{Linha do tempo de milestones em seis meses e doze sprints.}
\end{figure}

\begin{longtable}{@{}L{2.4cm}L{2.4cm}L{6.1cm}L{2.5cm}C{2.2cm}@{}}
  \caption{Cronograma complementar de entregas importantes}\\
  \toprule
  \textbf{Periodo} & \textbf{Sprint} & \textbf{Entrega Principal} & \textbf{Ambiente} & \textbf{Status} \\
  \midrule
  \endfirsthead
  \toprule
  \textbf{Periodo} & \textbf{Sprint} & \textbf{Entrega Principal} & \textbf{Ambiente} & \textbf{Status} \\
  \midrule
  \endhead
  Mes 1 & Sprints 1 e 2 & Configuracao do ambiente, arquitetura base e MVP da interface. & Desenvolvimento & Concluido \\
  Mes 2 & Sprints 3 e 4 & Liberacao do primeiro modulo core para ambiente de homologacao. & Homologacao & Concluido \\
  Mes 3 & Sprints 5 e 6 & Primeira release em producao, versao beta interna. & Producao/Beta & Concluido \\
  Mes 4 & Sprints 7 e 8 & Ajustes de feedback dos usuarios e refatoracao de performance. & Homologacao & Concluido \\
  Mes 5 & Sprints 9 e 10 & Implementacao de recursos avancados de seguranca e relatorios. & Homologacao/Producao & Concluido \\
  Mes 6 & Sprints 11 e 12 & Testes de carga, ajustes finais de UI/UX e rollout final. & Producao & Concluido \\
  \bottomrule
\end{longtable}

% ---------------------------------------------------------------------------
\section{Retrospectiva Consolidada -- Licoes Aprendidas}

\subsection{O que funcionou muito bem -- Sucessos}

\begin{successbox}{Sucessos observados no ciclo Scrum}
  \begin{itemize}[leftmargin=1.2em, itemsep=0.5em]
    \item \textbf{Alinhamento nas Dailies:} Comunicacao diaria eficiente reduziu impedimentos complexos.
    \item \textbf{Refinamento Tecnico -- Grooming:} As sprints da metade final do projeto foram melhor estimadas devido ao refinamento robusto.
    \item \textbf{Envolvimento dos Stakeholders:} Reviews quinzenais com feedbacks construtivos evitaram retrabalho no final.
    \item \textbf{Priorizacao do Backlog:} O time conseguiu adaptar o escopo conforme valor de negocio e urgencia tecnica.
  \end{itemize}
\end{successbox}

\subsection{Desafios Enfrentados e Como Foram Superados}

\begin{longtable}{@{}L{3.1cm}L{2.3cm}L{3.1cm}L{4.1cm}L{3.1cm}@{}}
  \caption{Desafios, solucoes e resultados}\\
  \toprule
  \textbf{Desafio} & \textbf{Periodo} & \textbf{Impacto} & \textbf{Solucao Aplicada} & \textbf{Resultado} \\
  \midrule
  \endfirsthead
  \toprule
  \textbf{Desafio} & \textbf{Periodo} & \textbf{Impacto} & \textbf{Solucao Aplicada} & \textbf{Resultado} \\
  \midrule
  \endhead
  Gargalo em Code Review & Sprints 3 e 4 & Atraso em entregas & Adocao do Buddy System e primeira hora do dia dedicada a revisoes & Reducao do tempo de revisao. \\
  Mudanca de Infraestrutura & Mes 4 & Instabilidade temporaria & Criacao de task force focada em DevOps & Pipeline de CI/CD estabilizado. \\
  Mudancas de Prioridade & Ao longo do projeto & Replanejamento do backlog & Refinamento continuo com PO e stakeholders & Escopo ajustado sem comprometer a entrega principal. \\
  \bottomrule
\end{longtable}

% ---------------------------------------------------------------------------
\section{Proximos Passos -- Proximo Ciclo}

Apos a entrega principal, o produto entra em fase de sustentacao, monitoramento e evolucao continua. Essa etapa deve garantir estabilidade operacional, coleta de feedbacks reais e priorizacao de melhorias para um possivel novo ciclo de desenvolvimento.

\begin{tcolorbox}[
  colback=white,
  colframe=PrimaryBlue,
  title=\textbf{Plano Pos-Entrega},
  colbacktitle=PrimaryBlue,
  coltitle=white
]
  \begin{minipage}[t]{0.31\textwidth}
    \begin{successbox}{Suporte e Sustentacao}
      Definir a equipe responsavel pelo monitoramento do ambiente produtivo, atendimento de chamados, acompanhamento de incidentes e manutencao corretiva/evolutiva.
    \end{successbox}
  \end{minipage}
  \hfill
  \begin{minipage}[t]{0.31\textwidth}
    \begin{infobox}{Backlog Remanescente}
      Consolidar itens despriorizados ao longo dos seis meses, mantendo-os organizados para uma possivel Fase 2 ou novo ciclo de evolucao do produto.
    \end{infobox}
  \end{minipage}
  \hfill
  \begin{minipage}[t]{0.31\textwidth}
    \begin{notebox}{Coleta de Feedbacks}
      Realizar pesquisa de NPS ou satisfacao com os usuarios finais nos primeiros 30 dias de uso da solucao, coletando dados para priorizacao de melhorias futuras.
    \end{notebox}
  \end{minipage}
\end{tcolorbox}

\begin{table}[H]
  \centering
  \caption{Acoes recomendadas para o proximo ciclo}
  \begin{tabularx}{\textwidth}{@{}Y L{3.2cm}L{2.4cm}C{2.4cm}@{}}
    \toprule
    \textbf{Acao Recomendada} & \textbf{Responsavel} & \textbf{Prazo} & \textbf{Prioridade} \\
    \midrule
    Monitorar ambiente produtivo e incidentes iniciais & \placeholder{Nome/Papel} & \placeholder{30 dias} & Alta \\
    Consolidar backlog remanescente para Fase 2 & \placeholder{Product Owner} & \placeholder{15 dias} & Media \\
    Realizar pesquisa de satisfacao/NPS & \placeholder{UX/Produto} & \placeholder{30 dias} & Alta \\
    Planejar roadmap de evolucao continua & \placeholder{Squad/PO} & \placeholder{45 dias} & Media \\
    \bottomrule
  \end{tabularx}
\end{table}

% ---------------------------------------------------------------------------
\section{Encerramento e Assinaturas}

Ao assinar este documento, as partes declaram-se cientes e de acordo com a entrega do escopo final acordado, reconhecendo o encerramento do ciclo principal deste projeto e a transicao para a etapa de sustentacao e evolucao continua.

\vspace{1.2cm}

\begin{table}[H]
  \centering
  \begin{tabularx}{\textwidth}{@{}C{0.31\textwidth}C{0.31\textwidth}C{0.31\textwidth}@{}}
    \rule{0.29\textwidth}{0.4pt} & \rule{0.29\textwidth}{0.4pt} & \rule{0.29\textwidth}{0.4pt} \\
    \placeholder{Nome do Product Owner} & \placeholder{Nome do Scrum Master} & \placeholder{Nome do Stakeholder Principal/Cliente} \\
    Product Owner & Scrum Master & Sponsor do Projeto \\
  \end{tabularx}
\end{table}

\vspace{1.2cm}

\begin{notebox}{Registro de Encerramento}
  \textbf{Data da assinatura:} \placeholder{DD/MM/AAAA}

  \textbf{Observacoes finais:} \placeholder{Inserir observacoes finais, restricoes, pendencias aceitas ou encaminhamentos formais.}
\end{notebox}

\end{document}
